\documentclass[aspectratio=169,10pt]{beamer}
\usetheme[numbering=fraction,progressbar=frametitle]{metropolis}
\usepackage{fontspec}
\setsansfont{Latin Modern Sans}
\usepackage{amsmath,amssymb,mathtools,graphicx,booktabs}
\usepackage{polyglossia}
\setdefaultlanguage{french}
\setbeamertemplate{navigation symbols}{}
\setbeamerfont{frametitle}{size=\large}
\setbeamerfont{normal text}{size=\normalsize}
\setlength{\parskip}{.45em}
\hypersetup{pdfauthor={},pdfsubject={Lecture slides},pdfkeywords={},pdftitle={Géométrie repérée : droites, projections et cercles}}
\providecommand{\N}{\mathbb{N}}
\providecommand{\Z}{\mathbb{Z}}
\providecommand{\D}{\mathbb{D}}
\providecommand{\Q}{\mathbb{Q}}
\providecommand{\R}{\mathbb{R}}
\providecommand{\C}{\mathbb{C}}
\providecommand{\vect}[1]{\overrightarrow{#1}}
\providecommand{\vi}{\vec{\imath}}
\providecommand{\vj}{\vec{\jmath}}
\providecommand{\vk}{\vec{k}}
\providecommand{\norm}[1]{\left\lVert #1\right\rVert}
\providecommand{\vc}[2]{\begin{pmatrix}#1\\#2\end{pmatrix}}
\providecommand{\deter}[2]{\det\left(#1,#2\right)}
\providecommand{\abs}[1]{\left\lvert #1\right\rvert}
\providecommand{\dd}{\,\mathrm{d}}
\providecommand{\e}{\mathrm{e}}
\providecommand{\Card}{\operatorname{Card}}
\providecommand{\E}{\operatorname{E}}
\providecommand{\V}{\operatorname{V}}
\providecommand{\degre}{^{\circ}}
\providecommand{\un}[1]{\,\mathrm{#1}}
\providecommand{\sci}[2]{#1\times10^{#2}}
\providecommand{\pgcd}{\operatorname{PGCD}}
\providecommand{\ppcm}{\operatorname{PPCM}}
\title{Géométrie repérée : droites, projections et cercles}
\subtitle{Cours et exemples}
\date{}
\author{}
\begin{document}
\begin{frame}[plain]
\vfill
{\usebeamerfont{title}\inserttitle\par}
\medskip
{\usebeamerfont{subtitle}\insertsubtitle\par}
\medskip
{\usebeamercolor[fg]{progress bar}\rule{\textwidth}{.5pt}}
\vfill
\end{frame}
\begin{frame}{Le parcours du chapitre}
\textbf{Objectifs.} Utiliser un vecteur normal ; calculer un projeté ; déterminer ou reconnaître un cercle ; traduire une configuration en équations.

\textbf{Prérequis.} Coordonnées vectorielles, produit scalaire, systèmes linéaires, complétion du carré. Repère toujours orthonormé.

\textbf{Repère annuel.} 5, 7 et 12 janvier : 6 h
\end{frame}
\begin{frame}{Vecteur normal et équation cartésienne}
Un vecteur normal non nul à une droite est orthogonal à tout vecteur directeur de cette droite. Si \(\vec n=(a;b)\ne(0;0)\) est normal et \(A(x_A;y_A)\) appartient à d, alors \(d\) a pour équation \(a(x-x_A)+b(y-y_A)=0\), soit \(ax+by+c=0\). Un vecteur directeur possible est \((-b;a)\).
\end{frame}
\begin{frame}{Vecteur normal et équation cartésienne}
Par \(A(1;2)\) et normale \((2;-3)\) : \(2(x-1)-3(y-2)=0\), donc \(2x-3y+4=0\). Le vecteur \((3;2)\) dirige la droite : produit \(6-6=0\).
\end{frame}
\begin{frame}{Vecteur normal et équation cartésienne}
\((a;b)\) est normal, et non directeur en général. \((a;b)\) ne doit pas être nul.
\end{frame}
\begin{frame}{Vecteur normal et équation cartésienne}
Équation de la droite par \((2;-1)\) de normale \((3;4)\).
\end{frame}
\begin{frame}{Projeté orthogonal et distance}
Le projeté H de P sur d est caractérisé par \(H\in d\) et \(\vect{PH}\) colinéaire à un vecteur normal de d. Pour \(d:ax+by+c=0\), écrire \(H=P-\lambda(a;b)\), avec \(\lambda =\dfrac{ax_P+by_P+c}{a^2+b^2}\). La distance est \(PH=\dfrac{|ax_P+by_P+c|}{\sqrt {a^2+b^2}}\). La formule de distance est un complément déduit, pas un automatisme isolé.
\end{frame}
\begin{frame}{Projeté orthogonal et distance}
\begin{columns}[T,onlytextwidth]
\begin{column}{.57\textwidth}
Pour \(P(3;4)\) et \(d:x+y-2=0\), \(\lambda =\dfrac{3+4-2}{2}=\dfrac{5}{2}\), donc \(H(\dfrac{1}{2};\dfrac{3}{2})\) et \(PH=\dfrac{5}{\sqrt 2}=\dfrac{5\sqrt 2}{2}\). Vérifier \(x_H+y_H=2\).
\end{column}
\begin{column}{.4\textwidth}
\centering
\includegraphics[width=\linewidth,height=4.5cm,keepaspectratio]{figures/premiere-07-1.pdf}
\end{column}
\end{columns}
\end{frame}
\begin{frame}{Projeté orthogonal et distance}
La distance est positive ou nulle : conserver la valeur absolue. Deux vérifications, appartenance et perpendicularité, sécurisent le projeté.
\end{frame}
\begin{frame}{Projeté orthogonal et distance}
Projeter \(P(4;1)\) sur \(d:x=2\).
\end{frame}
\begin{frame}{Équation d’un cercle}
Le cercle de centre \(I(\alpha;\beta)\) et de rayon \(r>0\) a pour équation \((x-\alpha)^2+(y-\beta)^2=r^2\). Une équation \(x^2+y^2+ux+vy+w=0\) se traite par complétion des deux carrés. Si le rayon au carré obtenu est négatif, l’ensemble est vide ; s’il est nul, c’est un point.
\end{frame}
\begin{frame}{Équation d’un cercle}
\(x^2+y^2-4x+6y+9=0\) devient \((x-2)^2+(y+3)^2=4\). Centre \((2;-3)\), rayon 2. Les signes des coordonnées du centre sont opposés à ceux figurant dans les parenthèses.
\end{frame}
\begin{frame}{Équation d’un cercle}
Une équation quadratique en x et y n’est pas toujours un cercle : vérifier les coefficients et le rayon réel.
\end{frame}
\begin{frame}{Équation d’un cercle}
Reconnaître \(x^2+y^2+2x-8y+17=0\).
\end{frame}
\begin{frame}{Résoudre une configuration}
Une intersection est obtenue en résolvant simultanément les équations. Pour un cercle et une droite parallèle à un axe, substituer la coordonnée constante donne une équation du second degré. Pour une médiatrice, écrire l’égalité des distances au carré ; les termes quadratiques s’éliminent.
\end{frame}
\begin{frame}{Résoudre une configuration}
Le cercle \(x^2+y^2=25\) et la droite \(x=3\) se coupent en \((3;4)\) et \((3;-4)\). La médiatrice de \(A(0;0)\) et \(B(4;2)\) vérifie \(x^2+y^2=(x-4)^2+(y-2)^2\), donc \(2x+y=5\).
\end{frame}
\begin{frame}{Résoudre une configuration}
Une intersection n’est pas l’union des ensembles de solutions. Vérifier les points dans les deux équations.
\end{frame}
\begin{frame}{Résoudre une configuration}
Intersecter \((x-1)^2+(y+2)^2=9\) avec \(y=1\).
\end{frame}
\end{document}
